\section{Exact reciprocity and the original irreducible-point normalization}
\label{sec:count-ii-producer}

\status{CIP1--CIP12 close two specified inputs of the original CountII
route.  They supply the exact reciprocity matching retained in JSP7,
and the complete \texttt{NiceI} instance of \texttt{Norm} used in
JSP11.  The latter statement is universal in the original point and
includes its integral representation, ordered characters, original
$D_0$, and determinant law.  The general \texttt{NiceE} instance is
outside this conclusion.  The ordinary arithmetic system, four further
Hida fields, all-level \texttt{HasAut}, and five other single-point
support inputs remain explicit.  In particular
\texttt{COUNT\_II\_SUPPORT} is not asserted to be closed.}

\subsection{The reciprocity convention on the actual local groups}

Retain the original CM field $L$, its maximal real subfield $L^+$,
the chosen algebraic closure $\Omega$, and the original
$s:\mathrm{HidaSetting}(\ell,L,\Omega)$.  For $v\mid\ell$ put
$w=s.w(v)$ and $K=L_w$.  The valuation on $\Omega$ is the given
$s.A(v)$, not a newly selected place.  KX3 identifies its decomposition,
inertia and Weil groups with those of $K$.  A Weil element of degree
one is an arithmetic Frobenius.  The source instead requires
\[
 \operatorname{ord}_K(a(\sigma))=-\deg(\sigma).
 \tag{CIP.1}
\]
Write $\operatorname{rec}^{\mathrm{ar}}_K$ for reciprocity sending a
uniformizer to arithmetic Frobenius.  Define its geometric counterpart by
\[
 \operatorname{rec}^{\mathrm{geo}}_K(u)
       =\operatorname{rec}^{\mathrm{ar}}_K(u)^{-1}.
\]
Thus the desired map from the Weil group is the inverse of geometric
reciprocity, not the inverse of arithmetic reciprocity.

The published input in this subsection is classical local and global
class field theory: the reciprocity isomorphisms, local--global and
norm compatibility, and the global kernel and surjectivity statements.
Their precise primary reference is J.~S.~Milne,
\emph{Class Field Theory}, version 4.03:
Introduction, p.~15; I.1.1 and I.1.10--1.11, pp.~20, 23--24;
III.3.2, p.~107; V.5.2--5.5 and Remark 5.7, pp.~178--180.
The inverse convention is imposed explicitly below.
\footnote{\url{https://www.jmilne.org/math/CourseNotes/CFT.pdf}.
The page numbers are the printed numbers; the PDF indices are larger
by eight.  V.4.10, pp.~175--176, also checks the inverse exponent on
rational local units.  The accompanying application certificate
separates these published inputs from the elementary matching proofs.}

% BEGIN CIP1
\paragraph{CIP1: the source-compatible global reciprocity map.}
\label{claim:CIP1}
For every number field $E$ and the specified algebraic closure
$\Omega$, inverse arithmetic reciprocity gives a map
\[
 \alpha_E:\mathbb A_E^\times\longrightarrow G_E^{\mathrm{ab}}
\]
satisfying all five fields of the source \texttt{IsGlobalArtin}.
\begin{proof}
Apply global reciprocity to $E$ and identify its abelian Galois group
with the topological abelianization of $\operatorname{Gal}(\Omega/E)$
by restriction to its maximal abelian subextension.  The theorem
applies to a number field, with its usual finite and infinite
completions; these are precisely the original idele factors.
Invert the arithmetic Artin map.  Inversion is a continuous
automorphism of this abelian target, so it preserves surjectivity and
the kernel
\[
 \ker\alpha_E=
 \overline{E^\times\, E_\infty^{\times,0}}.          \tag{CIP.2}
\]
Here $E_\infty^{\times,0}$ is exactly the source's
\texttt{infConn}: its real coordinates are positive and its complex
coordinates unrestricted.  This proves the \texttt{ker} field.

The continuous surjection is open.  For completeness, the idele group
is locally compact and a countable union of compact sets: at the
finite places the compact full-unit subgroup has countably many
valuation cosets, and the archimedean factors are countable unions
of compact annuli.  If $U$ is a neighborhood of one, choose a compact
neighborhood $C$ with $CC^{-1}\subset U$.  Countably many translates
of $C$ cover the domain.  Their images are closed compact sets
covering the compact Hausdorff target.  The Baire theorem gives
nonempty interior to one such image; consequently
$\alpha_E(C)\alpha_E(C)^{-1}$ contains a neighborhood of one.
It lies in $\alpha_E(U)$.  Applying this argument inside each
translate of an open set proves openness.  This gives the exact
\texttt{quot} field, including continuity and the quotient property.

Local--global compatibility identifies the image of local units with
the image of local inertia.  The unit group is compact, so its image
is closed, and the local reciprocity identification is onto this
inertia image.  This is the two-way \texttt{ine} assertion.  A
uniformizer under $\alpha_E$ has the image of an inverse arithmetic
Frobenius.  Multiplication by units runs over the whole inertia
coset, which is precisely the set of images of all such inverse
Frobenius elements.  This is the two-way \texttt{frob} assertion,
with the source's inverse sign.  Finally the closure of the local
multiplicative image is the image of the local decomposition group:
local reciprocity has dense image in its abelianization, and the
decomposition group is compact.  This is \texttt{dec}.  The local
identifications preserve the specified valuation on $\Omega$, as in
KX3, so none of these fields uses an unrelated local subgroup.
\end{proof}
% END CIP1

% BEGIN CIP2
\paragraph{CIP2: uniqueness under the literal global predicate.}
\label{claim:CIP2}
If $E$ is CM, every map satisfying the source
\texttt{IsGlobalArtin} equals $\alpha_E$ of CIP1.
\begin{proof}
Let $\alpha,\beta$ be two such maps, and let $q$ be any finite
continuous quotient of $G_E^{\mathrm{ab}}$.  The characters
$q\alpha,q\beta$ are continuous.  Choose a finite set $S$ of finite
places such that both kill the full product of unit groups away
from $S$.  This follows from an open neighborhood of one in their
kernels and the definition of the restricted-product topology.
Their ratio $\delta=(q\alpha)(q\beta)^{-1}$ kills diagonal $E^\times$
and all infinite ideles: all infinite places of a CM field are
complex and hence connected.

At a finite place $u\notin S$, choose a uniformizer $\pi_u$ and a
valuation over $u$ on $\Omega$.  The \texttt{frob} field for
$\alpha$ expresses $\alpha(j_u\pi_u)$ as the image of an inverse
arithmetic Frobenius $\sigma^{-1}$.  The reverse implication of the
\texttt{frob} field for $\beta$, using this same $\sigma$, supplies
another uniformizer $y_u$ with the same global image.  Since
$y_u/\pi_u$ is a unit killed by $q\beta$, the two maps agree at
$\pi_u$ after $q$.  They already agree on local units.  Therefore
$\delta$ is trivial on every $E_u^\times$ for $u\notin S$.

Fix $u\in S$ and $z\in E_u^\times$.  Weak approximation supplies
$a\in E^\times$ sufficiently close to $z$ at $u$ and to one at the
other places in $S$ that the corresponding local values of the
finite-valued continuous character $\delta$ agree.  On a diagonal
idele its value is one.  The factors away from $S$ and at infinity
are already trivial, so $\delta(j_u z)=1$.  Thus all local factors
are trivial.  An idele is a product of finitely many nonunit local
factors and an element of the full unit product; the latter is
handled by the same choice of $S$.  Hence $\delta=1$ on all ideles.
Finite continuous quotients separate points of the profinite group
$G_E^{\mathrm{ab}}$, giving $\alpha=\beta$ before taking quotients.
This argument uses the actual \texttt{frob}, \texttt{quot} and
\texttt{ker} fields, rather than assuming uniqueness as an extra
property of that predicate.
\end{proof}
% END CIP2

% BEGIN CIP3
\paragraph{CIP3: local-unit injectivity at the actual split CM places.}
\label{claim:CIP3}
Let $E$ be CM with complex conjugation $c$, and let $w\ne c(w)$
be a finite place.  If $u\in\mathcal O_{E_w}^{\times}$ and
$\alpha_E(j_wu)=1$, then $u=1$.  Every original $w=s.w(v)$ for
$v\mid\ell$ has this distinctness property.
\begin{proof}
Project (CIP.2) to the finite ideles.  The idele $j_wu$ belongs to
the closure of diagonal $E^\times$.  The full-unit subgroup
$U_f=\prod_z\mathcal O_{E_z}^{\times}$ is open, and
$E^\times\cap U_f=\mathcal O_E^\times$.  Therefore $j_wu$ lies in
the closure, inside $U_f$, of the diagonal global units.

For a global unit $e$, the element $e/c(e)$ is an algebraic integer
unit all of whose complex conjugates have absolute value one.
It is a root of unity.  One elementary proof is to take its powers:
their monic minimal polynomials have bounded degree and bounded
integer coefficients, since all their roots have modulus one.
There are only finitely many such polynomials and roots, so two
powers coincide.  The same argument bounds the set of all roots
of unity in the fixed field $E$, proving that $\mu(E)$ is finite.

The continuous map $y\mapsto y/c(y)$ on $U_f$ consequently carries
$j_wu$ into the finite, hence closed, diagonal copy of $\mu(E)$.
At any finite place outside $\{w,c(w)\}$ this idele has value one.
The embedding of $E$ in that completion is injective, so the root
of unity is one.  Thus $j_wu=c(j_wu)$.  Their $w$-coordinates are
$u$ and one because $w\ne c(w)$, proving $u=1$.

For the original data, $s.\mathrm{split}$ says that the
decomposition subgroup at $v$ in $G_{L^+}$ fixes $L$.
The image of this subgroup in $\operatorname{Gal}(L/L^+)$ is the
decomposition subgroup of $w/v$: surjectivity follows by extending
a place-preserving automorphism and adjusting its extension by an
element of $G_L$.  That image is trivial.  Since $L/L^+$ is
quadratic, $v$ has two distinct places $w,c(w)$ in $L$.
Equivalently, the completion in KX3 is $L_w\simeq L_v^+$.
Thus the hypothesis just used is supplied by the actual setting
at every place required by JSP7.  No assertion about nonsplit
places or a generic local-idele injectivity theorem is needed.
\end{proof}
% END CIP3

% BEGIN CIP4
\paragraph{CIP4: validity and uniqueness of the chosen \texttt{artOf}.}
\label{claim:CIP4}
For each original $v\mid\ell$, the source's
$\operatorname{artOf}(s.hw(v))$ satisfies \texttt{IsLocalArtin}.
Every map satisfying that predicate equals this map on the entire
original Weil group.
\begin{proof}
On $K=L_w$ use the local class field theory isomorphism
$K^\times\simeq W_K^{\mathrm{ab}}$, with its usual Weil topology,
and put
\[
 a_0(\sigma)
   =(\operatorname{rec}^{\mathrm{ar}}_K)^{-1}
            ([\sigma]^{-1})
   =(\operatorname{rec}^{\mathrm{geo}}_K)^{-1}([\sigma]).
 \tag{CIP.3}
\]
KX3's completion identification puts this map on the original
\texttt{Weil}(s.hw(v)).  It is onto, continuous on the original
inertia, and satisfies (CIP.1).  These give
\texttt{surj}, \texttt{cont}, and \texttt{ord\_eq}.

A smooth complex character of the Weil group has open kernel on
inertia, exactly as in the source definition.  Its kernel is open
in the Weil topology because inertia is open; the character is
therefore continuous.  It factors through the topological
abelianization, hence through $a_0$.  The quotient map to the
abelianization is open, so the resulting character of $K^\times$
has open kernel.  This proves the
literal \texttt{factor} field.  Local--global compatibility gives
\[
 \alpha_L(j_w a_0(\sigma))=\operatorname{toAb}(\sigma).
 \tag{CIP.4}
\]
CIP2 identifies every source \texttt{IsGlobalArtin} map with
$\alpha_L$, so (CIP.4) proves \texttt{glob} with its universal
quantifier, not just for one preferred global map.  Thus $a_0$
satisfies the full predicate.

If $a$ is any other map satisfying that predicate, then
$a(\sigma)/a_0(\sigma)$ is a unit by \texttt{ord\_eq}.
The \texttt{glob} field and (CIP.4) put its local idele in
$\ker\alpha_L$.  CIP3 forces this ratio to be one, for every
Weil element $\sigma$.  Thus $a=a_0$.  The source
\texttt{artOf} is the Hilbert-choice element of this nonempty
predicate, so it is valid and equals $a_0$ too.  In particular the
universal quantifier over local Artin maps in \texttt{IsWtPart}
can be replaced by this one map with no change of meaning.
Finally (CIP.1) and surjectivity show that the restriction from
inertia is onto $\mathcal O_K^\times$: a preimage of a unit has
degree zero, which is exactly inertia.
\end{proof}
% END CIP4

% BEGIN CIP5
\paragraph{CIP5: the positive cyclotomic norm and the exact JSP7 input.}
\label{claim:CIP5}
For every original inertia element $\sigma$ and every source
\texttt{IsLocalArtin} map $a$, one has, in
$\overline{\mathbb Q}_\ell$,
\[
 \epsilon(\sigma)=N_{K/\mathbb Q_\ell}(a(\sigma))
       =\prod_{\tau\in\mathrm{Emb}(\ell,w)}\tau(a(\sigma)).
 \tag{CIP.5}
\]
This supplies all of \texttt{JSP\_EXACT\_RECIPROCITY\_MATCH}.
\begin{proof}
First fix the sign over $\mathbb Q_\ell$.  Arithmetic reciprocity
on a unit $z\in\mathbb Z_\ell^\times$ acts on the
$\ell^m$th roots of unity by $z^{-1}$ for every $m$.
Here is a direct check consistent with the published example.
Choose a positive integer $b$, prime to $\ell$, sufficiently close
to $z^{-1}$ that $zb$ acts trivially on these roots under local
reciprocity, and also $b\equiv z^{-1}\pmod{\ell^m}$.
Continuity permits this choice.  Global arithmetic reciprocity
kills the principal idele $b$, and its real positive component.
At a prime $p\ne\ell$, the cyclotomic extension is unramified and
arithmetic Frobenius acts as the $p$th power.  Indeed its roots of
unity lie in unramified extensions by Hensel lifting of the
separable polynomial $X^{\ell^m}-1$ over finite residue fields.
The product of these away-from-$\ell$ actions is raising to
$\prod_{p\ne\ell}p^{v_p(b)}=b$.  The action of the $\ell$-component
$b$ is consequently raising to $b^{-1}$, and that of $z$ is
raising to $b$, hence to $z^{-1}$ modulo $\ell^m$.

Norm compatibility for local reciprocity now says that
$\operatorname{rec}^{\mathrm{ar}}_K(u)$ acts on all
$\ell$-power roots by $N_{K/\mathbb Q_\ell}(u)^{-1}$ for
$u\in\mathcal O_K^\times$.  CIP4 identifies
$\sigma$ in the local abelianization with
$\operatorname{rec}^{\mathrm{geo}}_K(a(\sigma))$.
Inverting the arithmetic action gives the positive norm in
(CIP.5).  This proves the identity for every admissible $a$.

The chosen valuation lies over both $v$ and $w$.  Since $v\mid\ell$,
its restriction places $\ell$ in $w$, so $K=L_w$ is a finite
extension of $\mathbb Q_\ell$ as required for the norm comparison.
The source set $\mathrm{Emb}(\ell,w)$ consists of continuous ring
embeddings $K\to\overline{\mathbb Q}_\ell$.  Such a map fixes
$\mathbb Q_\ell$ because it fixes $\mathbb Q$ and is continuous.
Conversely every $\mathbb Q_\ell$-embedding of the finite extension
$K$ is continuous by uniqueness of its extended valuation.
The extension is separable, so the norm is the product of exactly
these $[K:\mathbb Q_\ell]$ embeddings.  Thus the source's finite
product has precisely the factors in (CIP.5).

The definition of \texttt{epsV} is the original cyclotomic character
followed by the structural coefficient maps.  The field
$V.h\epsilon$ makes the same identification for the level frame.
Consequently (CIP.5) gives exactly the cyclotomic factor in
\texttt{IsTheta} and \texttt{IsUniv}, including its exponent $i$.
It is also \texttt{IsWtPart} with weight $-1$ and finite character
one: its displayed exponent is $-(-1)=1$.  Together with CIP4 this
is every validity, every-map uniqueness, and normalization clause
retained in JSP7.  Its proof of the increasing-weight and finite-part
dictionary can therefore be used without that additional input.
Separation of actual good arithmetic fibers remains a distinct
antecedent of JSP7; the reciprocity calculation does not produce it.
\end{proof}
% END CIP5

\subsection{The complete normalization needed at an original NiceI point}

Fix an odd prime $\ell$, the original $F:\mathrm{GFrame}(\ell)$,
a local $\mathbb Z_\ell$-algebra $O$, $n>0$, a residual datum $r$,
and an actual $x:\mathrm{Point}(r)$ satisfying \texttt{NiceI}.
Put $A=x.K[[t]]$, $F_A=x.K((t))$, and
$G'=\ker(\epsilon|_{G_0})$.  Here $G'$ uses the full
$\mathbb Z_\ell$-valued cyclotomic character of the frame.
The original representation and its form are
\[
 \rho:G_0\longrightarrow\mathrm{GL}_n(A),\qquad
 P=P^t\in\mathrm{GL}_n(A),\qquad
 \rho(g)^tP\rho(cgc^{-1})=\epsilon(g)^{1-n}P.       \tag{CIP.6}
\]
The invertibility of $P$ over $A$, not just $F_A$, is part of
the original \texttt{Point}.  Its ordered local characters are
$\psi_{v,0},\ldots,\psi_{v,n-1}$, and its flags are over $F_A$.

% BEGIN CIP6
\paragraph{CIP6: a finite constant extension standardizes the integral form.}
\label{claim:CIP6}
There is a finite extension $k'/x.K$, a coefficientwise map
$j:A\to A'=k'[[t]]$, and $T\in\mathrm{GL}_n(A')$ such that
\[
  T^t j(P)T=1.                                      \tag{CIP.7}
\]
The map $j$ is injective, local and finite free, sends $t$ to $t$,
and commutes with the original $\mathbb Z_\ell$-structure maps.
\begin{proof}
The composite $\mathbb Z_\ell\to O\to x.K$ makes $x.K$ a finite
field of characteristic $\ell$.  If its characteristic were a
different prime $p$, the unit $p\in\mathbb Z_\ell$ would map to
zero, which is impossible.  Moreover every map from
$\mathbb Z_\ell$ to this characteristic-$\ell$ ring factors
through the prime field $\mathbb F_\ell$; this also proves the
asserted compatibility after extension.  In particular two is a
unit.

Diagonalize $P$ integrally.  Modulo $t$, some vector has nonzero
self-pairing: otherwise the polarization identity
$2B(u,v)=B(u+v,u+v)-B(u,u)-B(v,v)$ would make the nondegenerate
bilinear form zero.  Lift this vector to a primitive vector in
$A^n$; its pairing with itself is a unit.  A primitive vector
extends to a basis over the local ring, and subtracting its
orthogonal projections from the other basis vectors gives a direct
orthogonal complement.  The complementary Gram matrix is again
invertible.  Induction gives
\[
 U^tPU=\operatorname{diag}(d_1,\ldots,d_n),\qquad
 U\in\mathrm{GL}_n(A),\quad d_i\in A^\times.
\]
Choose a finite extension $k'/x.K$ containing square roots of the
finitely many nonzero constant coefficients $d_i(0)$; the splitting
field of their quadratic polynomials is a finite field extension.
Each $j(d_i)$ has a square root $b_i\in A'^\times$.  Explicitly,
after choosing $b_i(0)$, the coefficient of $t^m$ in
$b_i^2=j(d_i)$ determines the new coefficient uniquely by division
by $2b_i(0)$.  Set $T=j(U)\operatorname{diag}(b_i^{-1})$.
This proves (CIP.7) with both $T$ and $T^{-1}$ integral.

Coefficientwise extension is injective and preserves the constant
coefficient, so it preserves and reflects units; this is the
required local-homomorphism property.  If $e_1,\ldots,e_d$ is an
$x.K$-basis of $k'$, decomposing every coefficient of a power series
in this basis gives a unique expression
$\sum_{a=1}^d e_a f_a(t)$ with $f_a\in A$.
Hence $A'$ is free of rank $d$ over $A$.  This proves ring-map
finiteness, rather than only finiteness of the residue extension.
\end{proof}
% END CIP6

% BEGIN CIP7
\paragraph{CIP7: the full block and its integral span.}
\label{claim:CIP7}
Define
\[
 \rho'(g)=T^{-1}j(\rho(g))T,\qquad
 \psi'_{v,i}=j\psi_{v,i}.                            \tag{CIP.8}
\]
The $F_A'=k'((t))$-span of $\rho'(G')$ is
$M_n(F_A')$.  Consequently for every $Z\in M_n(A')$ there is
$s\ge0$ such that
\[
 t^s Z\in\operatorname{span}_{A'}\{\rho'(g):g\in G'\}.
 \tag{CIP.9}
\]
Thus the source's block and span fields hold with $a=n$.
\begin{proof}
The original \texttt{TypeI} says exactly that $\rho|_{G'}$ is
irreducible after extension to an algebraic closure of $F_A$.
We give the finite-dimensional density argument it requires.
Over an algebraically closed field $C$, let $R$ be the $C$-span
of an irreducible group action on $V=C^n$.  It is a unital algebra,
since products of group matrices are group matrices.  Every
$R$-linear endomorphism of $V$ is scalar: it has an eigenvalue,
and the kernel of its difference with that scalar is a nonzero
invariant subspace, hence all of $V$.

For linearly independent $v_1,\ldots,v_r$, prove by induction that
$R\to V^r$, $a\mapsto(av_1,\ldots,av_r)$, is onto.  For $r=1$
this is irreducibility.  At the induction step its image $W$ is
an $R$-submodule whose projection on the first $r-1$ factors is
onto.  The kernel in the last factor is either zero or $V$.
In the latter case $W=V^r$.  In the former case $W$ is the graph
of an $R$-linear map $V^{r-1}\to V$.  Its component maps are scalar
by the preceding paragraph.  Applying this graph identity to
$a=1$ would express $v_r$ as a linear combination of its
predecessors, a contradiction.  At $r=n$ it follows that $R$
contains every endomorphism of $V$.

Apply this to the given absolute irreducibility.  The span over
$F_A$ must already have dimension $n^2$, since extension of scalars
preserves the dimension of a subspace and its matrix inclusion.
Thus finitely many matrices $\rho(g)$ form an $F_A$-basis of
$M_n(F_A)$.  After the coefficient extension and conjugation by
$T$ they still form a basis, proving the claimed full span.
Express a given integral $Z$ in this basis.  Each of its finitely
many coefficients lies in $k'((t))$; multiplying by a sufficiently
large power of $t$ makes all coefficients integral and proves
(CIP.9).  For $a=n$ every index of $\mathrm{Fin}(n)$ belongs to
the first block, so the literal space
$\mathrm{Ev}(A',A',n,n)$ is all of $M_n(A')$.
No factorial, trace denominator or characteristic-zero argument
has been used.
\end{proof}
% END CIP7

% BEGIN CIP8
\paragraph{CIP8: the original characters, torsion and integral flags.}
\label{claim:CIP8}
For every $v$, the characters in (CIP.8) are trivial on the
original subgroup $F.\mu(v)$ and admit an integral flag basis
$w_v\in\mathrm{GL}_n(A')$ in their original order:
\[
 w_v^{-1}\rho'(\phi_v(g))w_v
 \text{ is upper triangular with diagonal }
 (\psi'_{v,0}(g),\ldots,\psi'_{v,n-1}(g)).           \tag{CIP.10}
\]
Their complete original \texttt{N1} inequalities are preserved.
\begin{proof}
First prove the torsion assertion for the original $\psi_{v,i}$.
In the original rational flag basis, its value is a diagonal
entry of a conjugate of $\rho(\phi_v(g))$; it is therefore a
fixed finite $F_A$-linear combination of these matrix entries.
Clear all denominators by a power $t^N$.  The equation
$\psi_{v,i}(g)=1$ is then a finite integral linear equation in
these entries.  Modulo $t^m$ its solution set is closed, indeed
clopen, by the original entrywise \texttt{AdicCont} and
continuity of $\phi_v$.  Intersecting over $m$, and using
$\bigcap_m t^mA=0$, shows that $\ker\psi_{v,i}$ is closed.
It contains the commutator subgroup because the target is abelian,
so it contains its topological closure.

By definition $F.\mu(v)$ is the inverse image of the
$\ell$-primary torsion in the topological abelianization of $G_v$.
For $g$ in this subgroup, some $g^{\ell^b}$ belongs to that closed
commutator subgroup.  Therefore $\psi_{v,i}(g)^{\ell^b}=1$.
In the domain $A$ of characteristic $\ell$ this implies
$(\psi_{v,i}(g)-1)^{\ell^b}=0$, hence
$\psi_{v,i}(g)=1$.  Applying $j$ proves the same assertion for
$\psi'_{v,i}$.  This proof does not assume an additional continuity
field for the original characters.

Extend the given rational flag coefficientwise and conjugate it
by $T^{-1}$.  It is a full flag
$0=V_0\subset\cdots\subset V_n=(F_A')^n$ for $\rho'$ with
successive characters $\psi'_{v,0},\ldots,\psi'_{v,n-1}$.
Set $L_i=V_i\cap(A')^n$.  It is a finite free module of rank $i$:
it is a submodule of the finite free module over the DVR, and
clearing denominators of a basis of $V_i$ gives its rational rank.
It is saturated, since $az\in V_i$ for nonzero $a$ implies
$z\in V_i$.  The successive quotients $L_i/L_{i-1}$ are therefore
finite torsion-free modules of rank one, hence free.  Choose their
bases successively and lift them.  As $L_n=(A')^n$, the resulting
adapted basis is integral and invertible; let its matrix be $w_v$.
The action and its inverse preserve each $L_i$.  Its quotient
character agrees over the fraction field with the original ordered
$\psi'_{v,i}$, hence agrees integrally.  This proves (CIP.10).
Choose $\mathrm{ty}_v(j)=\mathrm{true}$ for every index.  Since
$a=n$, the antecedent in the original \texttt{w\_blk} field is
always false, so that field holds as well.

Finally injectivity of $j$ preserves inequalities of characters.
CIP6 identifies the two coefficient images of $\epsilon$.
An equality $\psi'_{v,i}=\psi'_{v,j}\epsilon\circ\phi_v$
would therefore descend to the forbidden original equality.
This proves both parts of \texttt{N1}, with exactly the original
index order $i<j$ in the second part.
\end{proof}
% END CIP8

% BEGIN CIP9
\paragraph{CIP9: the full determinant law and both nondegeneracy fields.}
\label{claim:CIP9}
The representation (CIP.8) has the image under $j$ of the entire
original determinant law, satisfies the source's normalized
\texttt{N3}, and satisfies \texttt{N4} for the original subgroup
$F.D_0$.
\begin{proof}
Let $C$ be any commutative $A'$-algebra and let
$y=\sum_g a_g[g]\in C[G_0]$.  Naturality of matrix extension and
(CIP.8) give
\[
 \sum_g a_g\rho'_C(g)
       =T_C^{-1}\left(\sum_g a_g\rho_C(g)\right)T_C.
\]
Determinants are unchanged by this conjugation.  This proves
\[
 \mathrm{GDet.ofRep}(\rho')
             =j_*\mathrm{GDet.ofRep}(\rho)           \tag{CIP.11}
\]
as polynomial laws on every coefficient algebra, rather than only
as characteristic polynomials of individual group elements.

For \texttt{N3}, take any original open set $U$ meeting $G'$.
The given $x.N3$ supplies $\sigma\in U\cap G'$ whose characteristic
polynomial over $F_A$ is separable.  After the field extension and
conjugation its characteristic polynomial is the coefficient image
of the same polynomial.  A Bezout identity with its derivative
survives the field extension, so it remains separable.  This is
exactly the normalized \texttt{N3} field, on the same $U$ and
the same $\sigma$.

For \texttt{N4}, over $F_A$ the original condition is that the
linear equations
\[
 \epsilon(g)\rho(g)X-X\rho(g)=0
                      \qquad(g\in F.D_0)            \tag{CIP.12}
\]
have only the zero solution in $n^2$ matrix entries.  Their
coefficient functionals consequently span the full dual of
$M_n(F_A)$; otherwise a nonzero annihilator would solve all of
them.  Choose a finite basis of these functionals.  The resulting
square coefficient matrix has nonzero determinant, which remains
nonzero over $F_A'$.  Thus (CIP.12) has no nonzero solution even
over the extended field.  If an integral matrix $Z\in M_n(A')$
satisfies the normalized equations, conjugate it by $T$ and view
it in $M_n(F_A')$.  By CIP6 the cyclotomic scalars agree, so
$TZT^{-1}$ satisfies the extended (CIP.12).  It is zero; hence
$Z=0$.  Throughout this argument the subgroup is the original
$F.D_0$.  It has not been replaced by the bottom subgroup stored
in \texttt{reIu}.
\end{proof}
% END CIP9

% BEGIN CIP10
\paragraph{CIP10: all fields of the exact NiceI normalization.}
\label{claim:CIP10}
For every original $I_u\le r.I_p(u)$ there exist
\[
 x':\mathrm{NP}(\mathrm{reIu}(F,I),n),\qquad
 j:x.A\longrightarrow x'.A
\]
with $j$ injective, local and finite, satisfying
\[
 \begin{gathered}
 \mathrm{GDet.ofRep}(\mathrm{mat}\,x'.\rho)
       =(\mathrm{GDet.ofRep}(\mathrm{mat}\,x.\rho)).\mathrm{map}(j),
 \qquad x'.\psi_{v,i}=\mathrm{pushCh}(j,x.\psi_{v,i}),\\
 x'.N3,\qquad x'.N4(F.D_0),\qquad x'.a=n.
 \end{gathered}                                      \tag{CIP.13}
\]
These are the complete conclusions of the original \texttt{Norm}
contract when its point hypothesis is \texttt{NiceI}.
\begin{proof}
Take $k',j,T$ from CIP6 and the representation and characters of
(CIP.8).  Give $k'$ the composite
$\mathbb Z_\ell\to O\to x.K\to k'$ algebra structure.
The coefficient field is finite.  Every entry of $\rho'$ is a
finite $A'$-linear combination of coefficient images of entries
of $\rho$ with fixed integral coefficients from $T,T^{-1}$.
The map $j$ preserves the $t$-adic filtration.  Reducing modulo
$t^m$ therefore proves the exact entrywise \texttt{AdicCont}
required for every $m$.

Applying $j$ to (CIP.6) and conjugating by $T$ gives
\[
 \rho'(g)^t\rho'(cgc^{-1})=\epsilon(g)^{1-n}1.
\]
Indeed $T^t j(P)T=1$ implies
$(T^{-1})^tT^{-1}=j(P)$; substituting this in the left side
gives the displayed identity.  CIP6 gives the exact
$\mathbb Z_\ell$-coefficient convention.  Thus \texttt{form}
has the literal matrix one.  Set $a=n$; $1\le a\le n$
follows from $n>0$, and CIP7 gives \texttt{blk} and \texttt{span}.

If $g\in I_u\le r.I_p(u)$, the original point has $\rho(g)=1$.
Coefficient extension and conjugation preserve this equality,
giving \texttt{ine} for the exact new tame subgroup.
CIP8 provides \texttt{psi\_mu}, \texttt{w}, \texttt{ty},
\texttt{w\_blk}, \texttt{tri}, and both \texttt{N1} clauses.
This has checked every field in the original \texttt{NP} record.
The changes $D_0=I_0=\bot$ inside \texttt{reIu} do not change
$\Gamma,G_0,G',\epsilon,G_v,\phi_v$ or $\mu_v$.
CIP9 gives (CIP.11), \texttt{N3}, and the separately requested
\texttt{N4} for the old $D_0$.  The character identities are the
definition (CIP.8).  CIP6 gives injectivity, locality and finiteness
of the one map $j$.  Finally $a=n$ holds outright, and therefore
also under the last implication $x.\mathrm{TypeI}\Rightarrow a=n$
printed in \texttt{Norm}.

This constructs the required object; it does not invoke the source
theorem named \texttt{norm}.  It proves the universal restriction
with the original \texttt{NiceI} hypothesis.  The disjunctive
\texttt{NiceI}\,$\vee$\,\texttt{NiceE} statement as a whole is
not needed for the following consumer and is not asserted here.
\end{proof}
% END CIP10

% BEGIN CIP11
\paragraph{CIP11: the original h=0 consumer keeps its point kernel and ideal.}
\label{claim:CIP11}
CIP10 supplies the entire \texttt{JSP\_SINGLE\_POINT\_NORM}
obligation in JSP11.  On the original rank-four route it has
$x'.a=4$, is valid at $\ell=3$, and introduces no change to
the selected ideal $J$ or to the kernel of an original point map.
\begin{proof}
Fix an arbitrary point $x$ and map $f:B\to x.A$ in the literal
\texttt{P1core} quantifier.  By definition $x$ is \texttt{NiceI}.
The $h=0$ instance of the original support chain uses the singleton
family $\mathrm{Fin}(1)$ with this very point; it never calls
\texttt{Norm} at a \texttt{NiceE} point.  The exact subgroup
identification in JSP11 is
\[
 I_u=(\mathrm{globOf}(r,V).\mathrm{frame}).I_u
                         =r.I_p(u).
\]
Thus CIP10 applies, with the original frame and its original $D_0$.
It returns precisely the normalized point required before the
remaining one-point \texttt{Choice} call, including both
nondegeneracy fields, all flags and all characters.

Let $O$ act on $k'$ through $O\to x.K\to k'$.
Coefficientwise construction in CIP6 gives
$j\circ(O\to A)=O\to A'$ and
\[
 \operatorname{const}_{k'}\circ j
              =(x.K\hookrightarrow k')\circ\operatorname{const}_{x.K}.
 \tag{CIP.14}
\]
Consequently $jf$ is an $O$-algebra map and retains the actual
residual determinant and ordered residual characters after this
one constant-field embedding.  Equation (CIP.11) is the full
determinant matching required by the original pair.  Injectivity
gives
\[
 \ker(jf)=\ker f,
 \qquad J\subseteq\ker f\ \Longleftrightarrow\ J\subseteq\ker(jf).
 \tag{CIP.15}
\]
No new ideal has been chosen.  A further injective extension supplied
by \texttt{Choice} preserves this equality and the displayed laws
by JSP10.  The original $N4$ transport used there applies to the
$N4$ just proved in CIP9.

The field extension and Gram--Schmidt argument only require that
two is invertible.  The full matrix span and determinant arguments
work in every characteristic.  Thus all conclusions apply to the
original $n=4,\ell=3$ case.  The construction depends on the point
$x$, as the source existential allows; it neither enlarges the
external coefficient object $c'$ nor alters the already selected
$B,\Lambda,J$.
\end{proof}
% END CIP11

% BEGIN CIP12
\paragraph{CIP12: the reduced joint-support frontier and its quantifiers.}
\label{claim:CIP12}
In the sufficient route of JSP12, replace the retained exact
reciprocity matching by CIP1--CIP5, and the retained one-point
\texttt{Norm} input by CIP6--CIP11.  All other producer inputs
remain on their original objects.  This reduction applies to both
actual rank-four Links, with the original fixed-$\iota$ and common
coefficient-threshold order.
\begin{proof}
First fix the original tower and the identification
$\iota:\overline{\mathbb Q}_\ell\simeq\mathbb C$.
The selected joint-support route must then provide one coefficient
threshold, before all later $c'$, original Spaces, and their
universal data.  At each allowed $c'$ it uses both Links with the
same $\kappa$ and their distinct auxiliary primes.  CIP1--CIP5
apply to every original $s$ at its actual $\ell$-adic places,
independently of $\iota$ and without enlarging $c'$.
CIP6--CIP11 apply to every later \texttt{NiceI} point over every
actual $B$, using only a pointwise finite constant extension already
allowed by the original \texttt{Norm} conclusion.  These operations
do not exchange any quantifiers or add a threshold.

On each Link the still-required ordinary system produces its own
specific $R,H,M,\mathbb T$ and the one ideal
$J=\ker(B\to\mathbb T)=\operatorname{Ann}_B(M)$ of JSP5.
JSP6 and JSP7 then give \texttt{RingProps} and \texttt{Dense},
with the reciprocity antecedent now supplied.  Its four further
Hida endpoints must still hold on that exact $J$:
\texttt{PointTheoremBaseObject},
\texttt{arithmeticPointPurityDense},
\texttt{arithmeticPointClassicalityDense}, and \texttt{TwistOk}.
The original all-level \texttt{HasAut}(B,J) is also still needed.
JSP8 supplies only its already reconstructed base, not its higher
level modules or the common-exponent identities.

Of JSP11's six support inputs, the required first input is now
proved.  The remaining five are the literal one-point
\texttt{Choice}, the actual \texttt{ArPl}, the exact first-order
bound, the exact obstruction bound, and the one-point \texttt{Raw}
support theorem.  The original \texttt{CohP}, \texttt{LocCFT},
coefficient, level and nontrivial residual-cyclotomic side
conditions retain their exact scopes in JSP11.  Supplying these
remaining inputs yields \texttt{P1core} on the already selected
$J$ by the same proof, with CIP11 at its normalization step.
This is a sufficient conditional implication, not an assertion
that those producers have been supplied.

The two ideals belong to the two different original universal
rings; within each ring all Hida and support properties refer to
the same ideal.  Once the remaining data satisfy the common
threshold, JSP12, QX10 and YX2 still assemble the original CountII
conclusion.  For RH and two-primary target characters only the
separate CaseIII branch is vacuous by KA5--KA6; CountII, as well as
Cases I and V, remains substantive.  Thus the present closures
contribute to both RH and unrestricted GRH.  They do not resolve
the extra CaseIII arithmetic work needed for general characters.
\end{proof}
% END CIP12
